\chapter*{Abstract}
\addcontentsline{toc}{chapter}{Abstract}

Breast ultrasound is a useful adjunct to mammography, but its interpretation is operator dependent and image appearance is affected by speckle, acquisition settings, and lesion context. This dissertation develops a reproducible research prototype for binary classification of breast-ultrasound images using convolutional neural networks. The software integrates contrast-limited adaptive histogram equalisation (CLAHE), aspect-ratio-preserving square padding, transfer-learning backbones (EfficientNet-B0 and ResNet-50), and a custom convolutional baseline. It also provides a FastAPI research demonstration for controlled image inference and visual inspection.

The principal contribution is methodological rather than clinical: a single preprocessing implementation is shared by training, evaluation, command-line inference, and the web service; paired lesion masks are recognised under the naming conventions used by the supplied datasets; and evaluation emits machine-readable metrics, figures, and an audit trail. A read-only data audit records split sizes and flags cross-split subject identifiers before performance is interpreted.

The supplied local copies of BUSI, OASBUD, and BrEaST are not yet a valid subject-independent benchmark. The audit identifies subjects spanning OASBUD and BrEaST partitions and cannot verify patient-level separation in the renamed BUSI copy. An end-to-end run of the bundled EfficientNet-B0 checkpoint on the current 120-image test partition produced 78.33\% accuracy and an ROC-AUC of 0.8612; these values are reported solely as provisional engineering results because the cohort has not passed the split audit. Historical headline figures, pseudo-labelling ablations, and noise curves contained in legacy project assets are therefore not treated as empirical findings.

The dissertation sets out the corrected workflow required for a defensible final study: rebuild the cohort at subject level, preserve source identifiers and preprocessing manifests, retrain each model with fixed seeds, and generate every reported table and figure from held-out predictions. Until those steps are complete, the system remains an educational research prototype and is not a diagnostic or clinical decision-support device.
